\section{Additional Model Details}
\label{app:model_details}

\subsection{Explicit \(p\)-Fold Tucker Motivation}
\label{app:pfold_tucker_motivation}

The power-dictionary channel in \eqref{eq:polynomial_linked_tucker} can be
viewed as a tied version of an explicit high-order Tucker interaction. This
construction is useful because it separates two ideas: the high-order
interaction pattern that a polynomial channel represents, and the parameter
sharing that makes the proposed model compact.

For a Tucker backbone, a formal \(p\)-fold Tucker channel may be written as
\begin{equation}
  \widetilde{Y}^{(p)}_{\mathbf{i}}
  =
  \sum_{\mathbf{r}^{(1)},\ldots,\mathbf{r}^{(p)}}
  \left(
  \prod_{q=1}^{p}
  X_{r_1^{(q)},\ldots,r_N^{(q)}}
  \right)
  \prod_{n=1}^{N}
  \mathcal{A}^{(n,p)}_{i_n,r_n^{(1)},\ldots,r_n^{(p)}} .
  \label{eq:app_explicit_pfold_tucker}
\end{equation}
For \(p=1\), this is the ordinary Tucker entry formula. For \(p=2\), each
mode projection couples two latent Tucker index tuples, giving a quadratic
Tucker-style channel. More generally, \eqref{eq:app_explicit_pfold_tucker}
is the tensor analogue of augmenting a linear map with polynomial interactions
among latent coordinates.

The explicit construction is not the parameterization used by the estimator:
each mode projection \(\mathcal{A}^{(n,p)}\) has \(I_nR_n^p\) entries and a
separate set of such projections would be needed for every active degree.
Dictionary-linked Tucker keeps the same interaction pattern but ties each
\(p\)-fold projection to the original Tucker factor,
\begin{equation}
  \mathcal{A}^{(n,p)}_{i_n,r_n^{(1)},\ldots,r_n^{(p)}}
  =
  \prod_{q=1}^{p} A^{(n)}_{i_n,r_n^{(q)}} .
  \label{eq:app_shared_projection_tying}
\end{equation}
Substituting \eqref{eq:app_shared_projection_tying} into
\eqref{eq:app_explicit_pfold_tucker} gives exactly
\(\cS^{\od p}\), where
\(\cS=\tucker{\cX}{\bA^{(1)},\ldots,\bA^{(N)}}\). Thus a degree-\(p\)
power-link channel is a shared-projection \(p\)-fold Tucker interaction: it
couples \(p\) latent Tucker tuples while reusing the same mode factors.

This viewpoint explains why the model can behave like a rank-lifted tensor
representation without storing an untied high-order tensor block. The
additional trainable variables are the scalar link coefficients, while the
high-order features are induced by repeated use of the same backbone.
